%----------------------------------------------------------------------
\chapter{Referencial Teórico}
\label{chp:referencial}
%----------------------------------------------------------------------

\epigraph{\it ``Knowing how things work is the basis for appreciation, and is thus a source of civilized delight.''}{-- William Safire}

Esta seção deve conter a Revisão Bibliográfica da sua proposta de dissertação de acordo com as instruções normativas descritas.
\section{Referências}


As referências bibliográficas devem ser claras e uniformes.  Por exemplo: \citep{Knuth1984}, \citep{Abdalla2002} e \citep{Smith1999}.


\subsubsection{Exemplo de sub-seção}

\lipsum[1]

\section{Uso de Figuras}

A Figura~\ref{fig:exemplo} apresenta uma conjunto de servidores distribuídos entre vários países. Você deve sempre optar por utilizar figuras no formate \textbf{.eps}.

\begin{figure}[!htb]
	\centering	
    \includegraphics[width=0.4\textwidth]{fig/distributedDatabase.eps} % <- você deve utilizar apenas .eps
    \caption{Banco de dados distribuído.}
	\label{fig:exemplo}
\end{figure}


\section{Uso de Tabelas}

Exemplo de uso de tabela abaixo:

\begin{table}[ht]
    \centering
    \caption{Legenda da tabela.}
    \label{quad:labeldatabela1}
\begin{tabular}{c c c }
        \hline
        \textbf{Parâmetros} & \textbf{Valor 1} & \textbf{Valor 2}\\ \hline
        \textbf{A} & 10 & 20 \\ \hline
        \textbf{B} & 10 & 20 \\ \hline
        \textbf{C} & 10 & 20 \\ \hline
    \end{tabular}
\end{table}


\section{Exemplo de enumeração}

Exemplo de enumeração:

\begin{enumerate}
	\item alguma coisa;
	\item outra coisa;
	\item última coisa.
\end{enumerate}
